\begingroup
\footnotesize\linespread{0.90}\selectfont
\setlength{\parskip}{0pt}\tableofcontents
\endgroup
\clearpage
\chapter*{Mapa skryptu i oznaczenia}
\addcontentsline{toc}{chapter}{Mapa skryptu i oznaczenia}

\noindent\textbf{Jak przechodzić między częściami.}
Pierwsza część prowadzi od rozmaitości i algebry przez
homotopię, homologię, formy, metryki i krzywiznę.
Podstawy grafów i drzew potrzebne dalej są w
Sekcji~\ref{sec:grafy-drzewa}, a grupa Pin przy algebrze
Clifforda w Definicji~\ref{def:pin-grupa}.
Druga część skupia grupy Liego, geometrię Kleina,
teorię Cherna--Weila i Hodge'a.
Trzecia prowadzi przez transwersalność, teorię Morse'a,
kobordyzmy, chirurgię, klasyfikację powierzchni,
przypadek wymiaru pięć i sfery Milnora.
Czwarta łączy potok Ricciego z twierdzeniem Poincarégo
w wymiarze trzy i geometryzacją Thurstona; twierdzenia
Perelmana pozostają w niej wynikami zewnętrznymi.
Piąta opisuje formy przecięcia w wymiarze cztery,
uchwyty Cassona, twierdzenie Freedmana, przeszkodę
Donaldsona oraz egzotyczne struktury gładkie na $\R^4$.
Dodatki~\ref{app:wyniki-zewnetrzne},
\ref{app:homologia-singularna-euler} i
\ref{app:grupy-algebry-macierzowe} zbierają kolejno
wyniki zewnętrzne, rachunki homologiczne oraz
definicje klasycznych grup macierzowych i ich algebr.

\noindent\textbf{Ścieżka do chirurgii.}
\[
\begin{gathered}
\text{rozmaitości i wiązki}
\longrightarrow \text{homotopia i homologia}
\longrightarrow \text{przecięcia i klasy charakterystyczne}\\
\longrightarrow \text{Morse i uchwyty}
\longrightarrow h\text{- oraz }s\text{-kobordyzm}
\longrightarrow \text{chirurgia i grupy }L.
\end{gathered}
\]
Rozdział~\ref{ch:chern-weil} przygotowuje klasy
Pontriagina używane później przy sferach Milnora.

\noindent\textbf{Najczęściej używane oznaczenia.}
\begin{center}
\begin{tabular}{p{3.9cm}p{9.7cm}}
$M^n,\ \partial M$ & rozmaitość wymiaru $n$ i jej brzeg rozmaitościowy;\\
$T_pM,\ TM,\ \nu_M$ & przestrzeń styczna, wiązka styczna i stabilna wiązka normalna;\\
$\pi_i(X),\ H_i(X),\ H^i(X)$ & grupy homotopii, homologii i kohomologii;\\
$\partial,\ \dd$ & operator brzegu na łańcuchach oraz różniczka form i kochain;\\
$\chi(M),\ e(E),\ p_i(E)$ & charakterystyka Eulera, klasa Eulera i klasy Pontriagina;\\
$\Z[\pi],\ L_n^s(\Z[\pi])$ & pierścień grupowy oraz prosta grupa chirurgiczna;\\
$\mathcal N(X),\ \mathcal S^s(X)$ & normalne niezmienniki i prosty zbiór struktur.
\end{tabular}
\end{center}
Definicje i hipotezy obowiązują w oznaczonych
miejscach tekstu; ta tabela służy nawigacji.

\clearpage
